\histitem{弗雷德霍姆指标与扰动}

指标这一概念于1920年出现在弗里茨·诺特关于积分方程的工作中。在1904年国际数学家大会的演讲中，希尔伯特考察了黎曼在势理论中提出的一个问题：给定平面内一个有界开域 \(\Omega\) ，其边界为光滑简单闭曲线 \(\Gamma\) ，以及三个函数 \(a\) 、\(b\) 、\(f\) ，它们定义在 \(\Gamma\) 上，求一个函数 \(z: \overline{\Omega} \to \mathbf{C}\) ，它在 \(\Omega\) 上全纯、在 \(\overline{\Omega}\) 上连续，并满足

\[
a \mathcal {R} (z) + b \mathcal {I} (z) = f \quad \text { on } \Gamma .
\]

将 \(\Gamma\) 用实数 \(s \in [0, 2\pi]\) 参数化后，希尔伯特证明该问题归结为求解关于 \(\varphi = \mathcal{R}(z)\) 的一个“具有奇异核”的积分方程，即

\[
a (s) \varphi (s) + \int_ {0} ^ {2 \pi} \mathrm{K} (s, t) \varphi (t) d t = f (s).\tag{4}
\]

这里，核 \(\mathrm{K}(s, t)\) 对某个连续函数 A 具有形式 \(b(s) \cot\left(\frac{t-s}{2}\right) + \mathrm{A}(s, t)\) 。因此它沿对角线奇异，上述积分须按柯西主值的意义理解。

诺特 {[}53{]} 指出，与满足弗雷德霍姆择一定理的积分方程不同，对于上述形式的核 K 和非零右端项 \(f\) ，(4) 可能有非平凡解族。他证明解空间由整数

\[
n = \frac {1}{2 \pi} \int_ {\Gamma} d (\log (a - i b)).\tag{5}
\]

若 \(n < 0\) ，则不存在非平凡解；若 \(n \geqslant 0\) ，方程(4)有一个依赖于 2n 个参数的解族。诺特使用“指标”（Index）一词称呼整数 n，并在其中认出一个绕数；在复变函数研究中，这一概念自十九世纪末以来已属经典。

后来的工作继续研究需要扩展弗雷德霍姆理论的方程实例。不过，在弗雷德霍姆映射的概念被系统提炼出来之前，还要再过二十多年；与此同时，对紧算子的扰动的研究也日趋成熟。

1909年，H. 外尔 {[}81{]} 证明，若定义在 \(\ell_{\mathbf{C}}^{2}(\mathbf{N})\) 上的两个有界二次型之差是完全连续的，则它们的本质谱相同（参见 III，第89页，定理2）。此外，里斯在1916年的回忆录中（当然没有使用现代语言）注意到，巴拿赫空间 E 上的紧算子构成 \(\mathcal{L}(\mathrm{E})\) 的双边理想。然而，似乎直到20世纪40年代才有人沿着这两条评注所指示的道路探索。1941年，J. 卡尔金受 J. 冯·诺伊曼关于算子代数的工作启发（我们很快将有机会讨论，见第537页），指出在希尔伯特空间框架下研究模此理想的同余关系的意义 {[}10{]}。他说明了如今以自己名字命名的代数的结构（参见 III，第75页）如何使人能够简单地恢复外尔的结果。

与此同时，大约1942年，J. 迪厄多内 {[}15{]} 对研究赋范空间之间的连续线性映射 \(u\) 的扰动产生兴趣，所考虑的情形是扰动按算子范数的意义“很小”。他只研究 \(u\) 是严格态射且其核有限维或有限余维的情形；他证明 \(u\) 的任意“微小扰动”也具有同样性质，并说明对于弗雷德霍姆映射（参见 III，第40页的 \(cf. n^{\circ} 2\) ），指标局部恒定（III，第58页定理1），III，第55页§4呈现的若干结果也如此。

这两种思想在1950年前后汇合，此时 F. 阿特金森 {[}3{]}、I. 戈赫伯格 {[}26{]}、S. 米赫林 {[}46{]} 和 B. 尤德 {[}93{]} 研究紧算子的扰动。他们明确指出了它与诺特指标的联系，并提出弗雷德霍姆映射 {[}3，第8页{]} 和里斯映射 {[}93，第5节{]} 的概念。随后十年中，人们对后者展开了系统研究；本书第三章呈现的结果在这一时期大致取得了今天的形式。

这些理论通常在巴拿赫空间框架中展开；然而，各种应用促使人们将其推广到更广泛的拓扑向量空间类。因此，弗雷歇空间的情形于1954年自然出现在 H. 嘉当和 J-P. 塞尔的工作中 {[}13{]}，该工作证明紧复解析流形取值于相干层的上同调是有限维的，其证明调用了 L. 施瓦茨证明的变形结果 {[}66{]}（参见 III，第73页定理2）。

在公式(5)的情形，若注意到右端（绕数）在同伦下不变，指标在变形下的不变性便得到充分说明。这样的指标拓扑表达式很快因微分几何中产生的弗雷德霍姆应用而闻名。若 D 是紧 \(C^{\infty}\) 流形 X 上的微分算子（参见 VAR，§14），且 D 是“椭圆的”（参见{[}2，第1节{]}），则将 D 延拓到适当的 Sobolev 空间定义一个弗雷德霍姆映射。寻求一个利用 X 的“特征类”给出其指标的拓扑表达式的公式——这一研究由 I. 盖尔范德约1960年的工作等所提出 {[}23，第I卷，第65页{]}——将在1963年导向 M. 阿蒂亚和 I. 辛格的“指标定理” {[}2{]}。后者将引发分析、拓扑和微分几何交汇处的非凡发展。我们无法在此描述他们开辟的广阔领域；这些领域至今仍极其富有活力，并且不断从弗雷德霍姆理论汲取营养。

\histitem{部分算子与谱定理}

一般的希尔伯特空间谱理论在很大程度上是为解决量子力学的数学基础问题而创立的，尤其是在冯·诺伊曼的推动下。

关于量子理论直至 W. 海森堡那篇引入量子力学的奠基性论文 {[}28{]}的发展详述，参见 M. 贾默的工作 {[}35{]}。1924年夏发表该文与1929年发表文章 {[}50{]} 相距不到五年；在后者中，冯·诺伊曼以极为清晰严谨的方式呈现了希尔伯特空间上部分（常称“无界”）埃尔米特算子的全部基本结果。

通过给出希尔伯特空间的第一种公理化表述 {[}49{]}，冯·诺伊曼完成了一个酝酿已久的思想（参见 ÉHM，第267页），并阐明了施密特、弗雷歇、菲舍尔和里斯在1910年前已熟知的序列空间或函数空间之间的各种同构。这种抽象程度的提升使冯·诺伊曼在概念上取得重大进展。

部分算子的思想便是如此；它似乎仍以未成形的形态盘旋在希尔伯特之后重新研究弗雷德霍姆方程的几部著作之上（见上文第1号）。希尔伯特注意到，许多 Sturm--Liouville 型微分方程，特别是在无界区间情形，可化为弗雷德霍姆方程；但是要应用积分方法，必须放宽对核施加的条件，并离开弗雷德霍姆、希尔伯特和施密特的结果适用的框架。多年来，人们研究了越来越多的奇异核，直到开始考虑这样的方程：其左端只有在未知量属于 L²(I) 的子空间时才有意义。在这方面必须提到 E. 希尔布 {[}30{]} 1908年的工作、H. 外尔关于 Sturm--Liouville 理论的1910年工作 {[}82{]}，以及最重要的 T. 卡勒曼 {[}11{]} 1923年的工作，后者在非常一般的框架中展开。

通过统一这些问题，冯·诺伊曼使这些经典工作呈现全新面貌。他尤其强调对称算子与自伴算子之间的根本区别，并用抽象“边界条件”来解释对称算子的自伴扩张，从而推广了在 \(R^{n}\) 的有界开集上拉普拉斯算子的狄利克雷问题和诺伊曼问题中出现的众所周知的边界条件。

与此同时，冯·诺伊曼 {[}51{]} 提出了正规算子 \(^{(5)}\) 的概念，将其作为谱理论的自然框架，并强调由算子及其伴随生成的交换代数的作用。

依托这一精确的数学形式主义——它还揭示了海森堡与薛定谔两种观点的等价性——以及对实验步骤和结果的哥本哈根诠释，非相对论量子力学从此达到一种此后几乎未有变化的精致程度；尽管这套物理理论的结果惊人、甚至看似悖论，其有效性自那以后仍不断得到实验确认，而且精度非凡。

从形式观点看，谱理论的基础由冯·诺伊曼奠定。与此同时，在后者首次提出这些思想之后，M. 斯通于1928年至1930年间进行了类似研究，并获得密切相关的结果。斯通于1932年发表的已知结果的清晰完整表述 {[}72{]}，对希尔伯特谱理论的传播产生了重要影响。

然而，表述上的其他改进也对这些结果的传播很重要。事实上，冯·诺伊曼给出的谱定理（参见 IV，第266页定理1）长期以来被认为非常困难，主要是因为它是在向量值测度框架中表述的。

后来，P. 哈尔莫斯 {[}27{]} 强调通过 \(L^{2}\) 空间上的乘法算子解释谱定理，从而揭示了它在有界算子情形下本质上初等的一面。最近，S. 沃罗诺维奇（在希尔伯特模上的正规算子这一稍有不同的 C*-代数框架中 {[}92{]}）通过引入

“有界化”（参见 IV，第265页第2号），使得可以像处理自伴算子一样简单地处理正规部分算子。

此外，若 D 是 \(R^{n}\) 的开集 U 上的形式对称标量微分算子，则研究 D 在 \(L^{2}(U)\) 上的自伴扩张自然与研究 U 上的广义函数相联系（拉普拉斯算子的情形参见 IV，第242页第6号）。关于后者的历史，参见 J. 迪厄多内的著作 {[}16，第VII章{]}——这是一段曲折的历史，紧接在我们刚才述及的时期之后展开，其中 S. 索伯列夫于1936年 {[}69{]} 和 L. 施瓦茨十年后的决定性贡献值得一提 {[}65{]}。

从纯数学观点看，谱理论与黎曼几何以一种特别富有成果的方式相互作用。因此，M. 卡茨于1966年重新讨论 S. 博赫纳提出的问题，给出了一种引人注目且极其生动的表述 {[}36{]}：确定一个黎曼流形的哪些几何性质可以从拉普拉斯算子的谱性质推出。物理学家 A. 舒斯特于1882年已经预见到这个问题 {[}64{]}：要弄清振动系统发出的不同音调，在某些特殊情形下可能可以解决，也可能不行；但是若要解决逆问题，根据钟能够发出的声音来确定钟的形状，即便是最有技巧的数学家也会束手无策。（“确定振动系统发出的谐波，在某些特殊情形下可能可以解决，也可能不行；但如果要解决逆问题、从钟能够发出的声音确定其形状，最有技巧的数学家也会感到困惑。”）。

\histitem{谐波分析与群论的联系}

我们在《积分学》一书的历史注记中（参见 ÉHM，第275页）描述过，与热方程有关的问题如何引导傅里叶提出把任意函数表示为三角函数之和这一革命性思想。

我们无法在此追溯整个十九世纪傅里叶级数和积分的发展，也无法追溯该理论提出的问题对分析演变的深刻影响——其影响甚至颠覆了实数的观念，并促成集合论的诞生（参见 ÉHM，第42页）。但为了理解本书所阐述思想的形式，有必要说明谐波分析在1925年以后如何同群论和希尔伯特谱理论发生联系。

傅里叶思想与 R/Z 和 R 的群结构之间联系的认识，以及将这些思想推广到其他群，来得相当晚。

在有限群及其特征标理论的早期阶段，自然可以辨认出一些如今看来包含有限傅里叶分析雏形的结果。例如，Z/qZ 的可逆元群的对偶，以及表达一个元素的特征函数为特征标线性组合的正交关系，已隐约出现在1837年狄利克雷的文章 {[}17{]} 中；该文证明当 a 与 q 互素时，存在无穷多个素数 \(p \equiv a \mod q\) 。他借助高斯的思想；高斯在研究具有整数系数及固定判别式的二元二次型时引入了“特征标”一词（参见 {[}22，第230节{]}）。

但高斯和狄利克雷都没有使用群的语言。戴德金在1879年讲述这些工作时，才认识到它们潜在的作用，并为交换群定义了特征标的概念。H. 韦伯大大扩充了戴德金的评述：先是1882年 {[}77{]}，然后是1886年 {[}78，第112页{]}；他在那里引入有限阿贝尔群 A 的对偶群，并注意到后者同构于 A。他当时的目的是构造具有给定伽罗瓦群的域 Q 的阿贝尔扩张，并未考虑这类群上的谐波分析问题。

在非交换群情形，有限群表示理论的基础在20世纪初前后，继 G. 弗罗贝尼乌斯、W. 伯恩赛德和 I. 舒尔的工作之后得到坚实奠定（参见 ÉHM，第154页）。1905年以后，特征标的正交关系在理论组织中发挥关键作用已很清楚。但它们与谐波分析的亲缘关系仍隐藏不显。

我们应归功于赫尔曼·外尔将这些领域联系起来：他构想了紧群表示的一般理论。在1925年至1927年间发表的几部重要著作中，他认识到这类理论与傅里叶分析及希尔伯特谱理论之间的联系，并为许多后来的发现奠定基础。

是舒尔的一封信把外尔引向这条道路的，缘于二人对不变量理论的共同兴趣。我们在关于 Haar 测度的注记（ÉHM，第289–291页）中描述过，A. 胡尔维茨早在1897年就已构想利用“不变积分”来确定在正交群下不变的 \(R^{n}\) 上多项式。舒尔似乎直到1924年才知道胡尔维茨的结果；那时他突然理解到，不变积分可用于把自己为有限群建立的特征标理论和正交关系推广到群 \(\mathbf{O}(n)\)。外尔立即看出，舒尔的方法与埃利·嘉当的工作相结合，使得可以构造紧半单李群的不可约表示。外尔将嘉当和舒尔的思想融为一体的三篇论文发表于1925年 {[}83{]}，其中已经包含 LIE, IX, §6–7 的基本结果；参见 ÉHM，第328–330页。

但在次年一篇与学生 F. 彼得合写且同样著名的文章 {[}54{]} 中，外尔才揭示紧群表示论、傅里叶分析和希尔伯特谱理论之间的联系。这篇不足二十页的文章孕育了许多未来的发展。

彼得和外尔摆脱了对所研究群的结构所作的任何代数假设，只假定 G 是一个配备不变测度的紧拓扑群。对连通李群而言，这种测度的存在是显然的；不久之后 A. Haar 将在一般情形建立它，部分动机正来自此处呈现的结果（参见 ÉHM，第291页）。

他们研究的出发点是矩阵系数的正交性。舒尔注意到，对每个不可约表示等价类选取一个代表 \(\pi \in \widehat{\mathrm{G}}\) 、\(\pi\) 的空间 E 上一个 G-不变内积以及 E 的一个标准正交基，就得到一族在 \(\mathrm{L}^2 (\mathrm{G})\) 中正交归一的矩阵系数。彼得和外尔计划证明这族函数是完备的。

在有限群情形，弗罗贝尼乌斯曾通过研究群代数 C{[}G{]} 证明了相应结果。彼得和外尔接着考虑 G 上连续函数空间 \(\mathcal{C}(\mathrm{G})\) ，并为其赋予卷积乘法。这似乎是首次将卷积乘法作为抽象运算、明确地同群结构联系起来（参见 ÉHM，第295页）。此外，彼得和外尔把 \(\mathcal{C}(\mathrm{G})\) 的一个元素称为 Gruppenzahl（“群数”），并以 xy 表示两个 Gruppenzahlen 的（卷积）乘积。他们还为 \(\mathcal{C}(\mathrm{G})\) 配备对合 \(f \mapsto \widetilde{f}\) ，其中 \(\widetilde{f}(g) = \overline{f(g^{-1})}\) ，对于 \(g \in G\) 。

在引入对合代数 \(\mathcal{C}(\mathrm{G})\) 后，彼得和外尔的第一个观察是，每个 G 的酉表示 \(\pi\) 都给出一个表示 \(f \mapsto \pi(f)\) ，其代数为 \(\mathcal{C}(\mathrm{G})\) （参见 V，第400页）。他们明确给出了 \(\mathcal{C}(\mathrm{G})\) 的埃尔米特元素概念，也（隐含地）给出了正元素概念，为应用希尔伯特空间技巧打开道路。

彼得和外尔不再犹豫，称算子 \(\pi(f)\) 为 \(f\) 的“傅里叶系数”，并通过指出舒尔正交关系蕴含贝塞尔不等式，继续与傅里叶理论作平行比较

\[
\sum_ {\pi \in \widehat {\mathrm{G}}} \dim (\pi) \operatorname{Tr} (\pi (f) \pi (f) ^ {*}) \leqslant (f * \widetilde {f}) (e) = \| f \| _ {2} ^ {2}.\tag{6}
\]

希尔伯特--施密特谱理论随后使他们能够证明这是等式（“普朗谢雷尔公式”）。对 \(\varphi\) （\(\mathcal{C}(\mathrm{G})\) 中的每个元素），彼得和外尔关联连续核 \(K_{\varphi}\colon G\times G\to C\) ，其定义为 \(\mathrm{K}_{\varphi}(x,y)=\varphi(xy^{-1})\) 。随后他们观察到，如果 \(\varphi\) 形如 \(f*\widetilde{f}\) 且 \(f\in\mathcal{C}(\mathrm{G})\) ， \(^{(6)}\) 则 \(K_{\varphi}\) 是希尔伯特--施密特理论意义下的正埃尔米特核。施密特算法的一个变体 \(^{(7)}\) 使他们可以构造核 \(K_{\varphi}\) 的特征函数标准正交基，继而将(6)中等式的证明归结为：由 \(K_{\varphi}\) 定义的算子的迹等于其特征值之和。

当然，只有表示 \(\pi\) 满足 \(\pi(f) \neq 0\)才可能出现在(6)中；\(K_{\varphi}\) 的谱理论表明，只要 \(f\) 不恒为零，就存在这样的表示。彼得和外尔很容易由此推出，不可约表示分离 G 的点——这是盖尔范德--赖科夫定理的一个特例，稍后将讨论该定理。

在弗罗贝尼乌斯的著名工作中，与(6)类似的等式曾使人能够证明 G 的正则表示包含所有不可约表示。但这涉及将该等式用于 C{[}G{]} 的单位元 f；此时显然有 \(\pi(f) \neq 0\) 对所有 \(\pi\) 都成立 。依彼得和外尔之见，卷积没有单位元解释了他们证明中的许多困难；然而，他们借用了傅里叶级数理论中一个注定具有巨大影响的思想，指出可将(6)应用于一列构成卷积逼近单位的函数，从而推出弗罗贝尼乌斯结果的类似物（参见 I，第120页第10号）。

同样地，从(6)产生的等式出发，其极化形式

\[
\sum_ {\pi \in \widehat {\mathrm{G}}} \dim (\pi) \operatorname{Tr} (\pi (x) \pi (y)) = (x * y) (e)\tag{7}
\]

通过令 \(y\) 为卷积逼近单位，可得到 \(\mathcal{C}(\mathrm{G})\) 的每个元素都能由不可约表示的矩阵系数的线性组合一致逼近。这些方法刚刚使外尔能对 H. 玻尔的“几乎周期函数”理论的基本结果给出新证明 {[}84{]}；我们很快将回到该理论。彼得和外尔指出，其中可以看出对非紧群表示理论的第一个分析应用，这里该群是 \(\mathbf{R}\) 的平移群。

年轻的量子理论很快给外尔提供了重新回到 R 的表示的机会。1927年底，他注意到（无论有限与否的）群和表示对于阐明该理论基础的意义 {[}85{]}。次年，他在一本轰动性的著作 {[}86{]} 中再次讨论这一问题。

在1927年的文章中，外尔注意到，若 u 是希尔伯特空间上的连续自伴算子，则 \(t \mapsto e^{itu}\) 是 R 的酉表示。但存在并非这种形式的 R 的酉表示；重新采用冯·诺伊曼关于用对称部分算子表示可观测物理量的思想，外尔提出每个 R 的酉表示都形如 \(t \mapsto e^{itu}\) ，其中 u 是希尔伯特空间上的部分自伴算子。这一观点使他能够把海森堡和薛定谔的“正则关系”（冯·诺伊曼已作详细研究）的分析，归结为同一希尔伯特空间上两个 R 的酉表示之间的联系。斯通 {[}73{]} 在1932年、沿着他对自伴算子谱性质的研究，证明了外尔所期待的结果（V，第428页，定理1）。

\histitem{局部紧交换群}

到20世纪30年代初，谐波分析与群表示之间的互动已经确立。接下来的十年中，二者相互推动，互动大为加强，分别促成对方迅速发展。

这种互动的一个极其肥沃的土壤是几乎周期函数的研究；H. 玻尔于1924年在实变量函数的情形中引入了它们 {[}9{]}。根据他的一个基本定理，这些函数正是 \(\mathbf{R}\) 上可由 \(x \mapsto e^{i\lambda x}\) 、\(\lambda \in \mathbf{R}\) 的线性组合一致逼近的有界连续函数（参见 II，第292页，习题54）。博赫纳 {[}6{]} 于1926年证明，\(f \in \mathcal{C}_b(\mathbf{R})\) 是几乎周期的，当且仅当其平移 \(x \mapsto f(x - a)\) 对于 \(a \in \mathbf{R}\) 在一致收敛拓扑下构成 \(\mathcal{C}_b(\mathbf{R})\) 的相对紧子集。博赫纳和冯·诺伊曼注意到，若 G 是拓扑群，这一刻画显然在 G 上（右侧或左侧）提供了几乎周期函数的概念。20世纪30年代初，在冯·诺伊曼 {[}52{]} 及彼得和外尔的方法推动下，这一理论有望迅速发展。

另一个富有成果的领域出现在 N. 维纳工作的边缘。他在1932年一篇著名论文 {[}88{]} 中把代数方法引入所谓塔伯尔问题的研究——也就是在已知某些加权平均行为的情况下，研究函数（或序列）沿滤子的渐近行为。他的方法受到 R. 施密特的思想启发。

维纳工作中的一个辅助结果（{[}88，引理 IIe{]}）尤其令人印象深刻并激发了诸多发展：若 \(f\) 是一个傅里叶级数绝对收敛且 \(f\) 不为零的周期函数，则 \(1/f\) 的傅里叶级数绝对收敛（参见 I，第38页，例）。早在1932年，R. 佩利和 N. 维纳 {[}89{]} 就勾勒出这一收敛结果、离散交换群的对偶群概念与几乎周期函数理论之间的联系。

在这一分析背景下，代数拓扑促使 L. 庞特里亚金为局部紧阿贝尔群的特征标研究提供决定性推动。他早在1932年就宣布，特征标群概念的意义在于把欧氏空间紧子集的同调群与其补集的同调群置于对偶关系中 {[}55{]}。在1934年发表的论文 {[}58{]} 中，他证明离散群的特征标群是紧的，随后在这一框架中建立对偶定理。虽然庞特里亚金使用了彼得和外尔的结果以及刚刚在一般情形建立存在性的哈尔测度，但他的主要目标是拓扑应用 {[}57{]}，并未明确关注谐波分析 \(^{(8)}\)：这些对偶定理是通过细致研究紧阿贝尔群的结构证明的。

庞特里亚金的工作立即引起极大兴趣。应冯·诺伊曼的建议，E. 范坎彭次年将对偶性推广到所有局部紧交换群 {[}37{]}。不久之后，他把新不变分析应用到几乎周期函数的研究，给出一个引人注目的结果 {[}38{]}；A. 外尔在同一时期独立发现了该结果 {[}79{]}：若 G 是局部紧阿贝尔群，可以得到一个紧群 \(\widetilde{G}\)：先给对偶群 \(\widehat{G}\) 赋予离散拓扑，再把由此得到的离散群的紧对偶群定义为 \(\widetilde{G}\)。群 G 随后典范地嵌入 \(\widetilde{G}\)（参见 II，第292页，习题54），对 \(\widetilde{G}\) 上连续函数应用彼得--外尔理论的简单结果，就能得到当时关于 G 上几乎周期函数的所有已知结果。

与此同时，在这一时期，经典分析中早已熟知的正定函数出现在 R 的酉表示研究中。J. 默瑟在1910年以前不久、结合积分方程理论研究了普遍正核 {[}45{]}；而其序列类比则在同一时期引出大量著作，尤其与矩问题有关（参见 IV，第359页，习题21）。与此同时，M. 马蒂亚斯于1923年已系统研究 R 上的正定函数 {[}43{]}；在博赫纳手中，它们成为傅里叶分析和概率论的基本工具之一 {[}7{]}。1933年，博赫纳 {[}8{]} 与里斯 {[}61{]} 独立注意到，R 的酉表示的每个对角矩阵系数都是 R 上的正定函数。这一观察使他们能够给出斯通定理的更简单证明，并将对一般理论的发展产生决定性影响（参见第7号）。

这一时期的结果由 A. 外尔整理在一本1937年前完成、1940年出版的著作中 {[}80{]}。在那里，他详细阐述彼得--外尔、庞特里亚金和范坎彭的结果。但在交换群情形，庞特里亚金和范坎彭强调群的结构和对偶定理，而外尔系统地发展了谐波分析，引入傅里叶变换并证明普朗谢雷尔公式与博赫纳定理（V，第455页，定理5）。他尤其强调卷积（称为“合成乘积”）的作用，以及正定函数的作用；后者如今定义在任意群上，并同表示的对角矩阵系数联系起来，至少在有限维表示情形如此。

因此在交换情形达到的这种一般性，稍晚些时候将为数论开辟新的视野。早在1936年，为向类域论引入代数方法，C. 谢瓦莱 {[}14{]} 引入了数域的 idèle 群；该群后来会作为 adèle 环的可逆元群出现（参见 AC，VII，第221--222页及 ÉHM，第143页）。1950年，J. 泰特 {[}75{]} 将证明，对 adèle 和 idèle 群的谐波分析使得可以容易地恢复 Hecke L-函数的函数方程，预示了把 adèle 群上的谐波分析同自守形式研究联系起来的非凡发展。

可以说，到20世纪30年代末，不变谐波分析的主要结果已经在紧群和局部紧交换群上建立。在交换情形，外尔的证明仍像庞特里亚金和范坎彭的证明一样，依赖于对群结构的精细了解，粗略说即 II 第2节第244页的分类结果。下一个十年借助盖尔范德学派的新方法将说明如何摆脱这一依赖（参见 H. 嘉当和 R. 戈德芒 {[}12{]}）。

\histitem{算子代数}

我们已经看到，里斯在1913年阐释希尔伯特学派工作时，赋予希尔伯特空间自同态代数 \(\mathcal{L}(\mathrm{E})\) 的作用，也赋予全纯函数演算和谱的抽象概念以作用。在1916年的回忆录中，他似乎已经预示赋范代数是谱理论的自然框架；经过巴拿赫及其学派在20世纪20年代

的工作，这一思想占据中心地位也就不足为奇了。M. 长友 {[}48{]} 于1936年引入巴拿赫代数概念，而 S. 马祖尔 {[}44{]} 于1938年证明，唯一作为域的 C 上赋范代数就是 C 本身（I，第26页，推论2）。

大约在20世纪30年代中期，与谱理论相关，抽象考察希尔伯特空间上的算子代数的机会不断增多。F. 默里和 J. 冯·诺伊曼在一系列深刻且有影响的工作 {[}47{]} 中，系统研究了等于自身双交换子的 \(\mathcal{L}(\mathrm{E})\) 的对合子代数（“冯·诺伊曼代数”）。另一方面，S. 斯廷于1936年提出引入“算子代数”的公理化概念，并深入研究相关结构 {[}70{]}。此外，斯通 {[}74{]} 受谱投影研究启发，于1937年研究其所有元素均为幂等元的含单位代数（“布尔代数”：参见 ÉHM，第146页）。他注意到，若 X 是局部紧拓扑空间，则开集的特征函数和处处稠密集补集的特征函数在 \(\mathcal{C}(\mathrm{X})\) 中生成一个布尔代数 A。随后他证明 A 的理想自然对应于空间 X 的闭子集，从而引入了一个后来特别富有成果的思想。

从1939年起，盖尔范德的观点为这些相对分散的结果打开了新视野。他的工作似乎部分受到维纳思想的推动，特别是其关于绝对收敛傅里叶级数的定理，以及彼得和外尔的方法。1939至1946年间，他与 M. 奈马克、D. 拉伊科夫和 G. 希洛夫合作，奠定了交换巴拿赫代数理论和星代数理论的基础（参见 {[}23{]}，第I卷，第169--400页）。特别地，他给出了维纳定理的经典证明——由 P. 莱维 {[}39{]} 推广——本书中可见该证明（I，第38页，例）。

在盖尔范德的方法中，与交换巴拿赫代数 A 相关的本质对象是 A 的极大理想空间 X。事实上，通过盖尔范德--马祖尔定理（I，第26页，推论2），A 的每个元素都在 X 上定义一个函数。众所周知，这一观点后来会在代数几何的发展中变得十分重要（参见 ÉHM，第146--148页）。这一方法本身受到了上文提到的斯通关于布尔代数的工作的启发。

似乎是 L. 卢米斯在1953年出版的课程 {[}41，第53页{]} 中首次强调 A 的特征标空间来阐述盖尔范德理论，并按本书所理解的意义定义盖尔范德变换。这一观点也许最清楚的优点是，特征标空间的拓扑性质直接来自巴拿赫空间对偶单位球在弱拓扑下的紧性。

从最早的工作起，盖尔范德就通过柯西积分在巴拿赫代数中构造了一元全纯函数演算，并由此推出：当极大理想空间不连通时，对合巴拿赫代数允许按环的乘积作非平凡分解。任意巴拿赫代数的情形直到1953年才得到处理，那时希洛夫在多元情形发展出一种函数演算。{[}68{]}。

至于星代数，盖尔范德和奈马克早在1942年就已证明它们可以实现为希尔伯特空间上的算子代数（参见 V，第442页定理2），它们很快成为大量深刻研究的对象；推动力来自群表示的应用、默里和冯·诺伊曼的工作以及量子力学的数学形式化。特别值得一提的是 I. 西格尔的贡献；他是系统研究星代数表示的先驱之一，这项研究与局部紧群的酉对偶研究相联系，我们将在几行后讨论后者。

尽管盖尔范德、奈马克、拉伊科夫和希洛夫在1946年以前已经奠定 C*-代数理论的基本基础，某些技术细节后来才得到改进。例如，R. 阿伦斯 {[}1{]} 说明了如何避免调用“希洛夫边界”（参见 I，第171页，习题26）的存在性，以证明盖尔范德变换是一个对合态射。此外，盖尔范德--奈马克关于含单位 C*-代数的理论起初把 \(1 + x^{*}x\) 对所有 \(x\) 可逆这一事实作为公理加入。然而他们猜想这个公理并非必要；这一断言的证明虽然初等却很精细，直到1952年前后才在 M. 福卡米亚 {[}21{]} 的工作、经 I. 卡普兰斯基 {[}62{]} 接续之后给出。

\histitem{局部紧群的表示}

直到1939年，拓扑群表示的研究局限于紧情形和交换情形。另一方面，物理学家在量子力学中发现了许多理由，需要研究不可约酉表示，以寻找能够模拟基本粒子状态的希尔伯特空间。P. 狄拉克在 E. 马约拉纳工作的基础上，于1936年前后指出，在这一视角下，洛伦兹群 SO(3,1) 和庞加莱群 SO(3,1) × \(R^{4}\) 的意义。他们的方法仍与当时数学家的方法相去甚远。但1939年，E. 维格纳 {[}90{]} 以默里和冯·诺伊曼的工作为基础，对庞加莱群的不可约酉表示作出分类，打开了一道缺口。

我们似乎应将研究局部紧群表示的一般路径归功于盖尔范德。他在1942年写的一篇回忆录中（参见 {[}23{]}，第II卷，第3--17页）与拉伊科夫注意到，这一理论之所以可能，是因为酉表示与正定函数之间存在联系。他们证明每个正定函数都有希尔伯特实现（V，第432页定理1），并由此推出不可约酉表示分离 G 的点（V，第454页定理4）。对合代数 L \(^{1}\) (G) 上的正线性泛函发挥了基本作用：盖尔范德与拉伊科夫揭示了它们同正定函数相联系的关系（参见 V，第448页命题13）。对合代数的包络 C* 代数在 {[}24，第48节{]} 中已经（未命名）出现，不过当时还没有直接提出把代数 Stell(G) 与 G 联系起来（参见 I 第125页定义9）。

1945年以后，酉表示的研究突飞猛进。与 C*-代数的互动以及对量子理论基础的分析很快带来有价值的一般结果：约1947年，西格尔把各种 \(^{(9)}\) 星代数同一个局部紧群联系起来 {[}67{]}，并编织起它们的表示与所研究群的表示之间的联系。不久后，G. 麦基提炼出诱导表示的概念 {[}42{]}，它将在酉表示的具体结果中无处不在。

与此同时，对揭示其多样性与深度的例子的研究大大推动了该理论。盖尔范德和奈马克研究 SL(2, C) 这一例子，于1947年分离出诱导表示概念的一个特殊情形，随后证明它是研究该群表示的关键 {[}23，第II卷，第41--124页{]}。同年，V. 巴格曼研究群 SL(2, R)，发现一个可数的平方可积表示族，称为“离散系列”，以及它们的矩阵系数之间的正交关系（参见 V，第424页命题8） {[}5{]}。正如戈德芒立即认识到的 {[}25{]}，Bargmann 正交关系由所有局部紧群的平方可积表示满足。

很快，对特殊群类的研究又一次、并且长时期内优先于抽象理论。特别是，哈里什-钱德拉将着手一般研究约化李群的表示——这项史诗般的任务，以及朗兰兹将为数论预见的卓越后果，我们都无法在此描述。

\specialchapter{参考文献}

{[}1{]} R. ARENS - "关于盖尔范德和奈马克的一个定理", Proc. Natl. Acad. Sci. USA 32 (1946), p. 237-239, Zbl 0060.27101.

{[}2{]} M. F. ATIYAH and I. M. SINGER -- "紧流形上椭圆算子的指标", Bull. Amer. Math. Soc. 69 (1963), p. 422--433, Zbl 0118.31203; Michael Atiyah 文集, vol. 3, Clarendon Press, 1988, p. 13--24.

{[}3{]} F. V. ATKINSON -- "赋范空间中线性方程的正规可解性", Mat. Sb., Nov. Ser. 28 (1951), p. 3--14, Zbl 0042.12001（俄文）。

{[}4{]} S. BANACH - 线性运算理论, Warsaw, 1932, Zbl 0005.20901.

{[}5{]} V. BARGMANN - "洛伦兹群的不可约酉表示", Ann. of Math. (2) 48 (1947), p. 568-640, Zbl 0045.38801.

{[}6{]} S. BOCHNER -- "几乎周期函数理论的贡献。II：多变量函数", Math. Ann. 96 (1926), p. 383--409, JfM 52.0265.01; 论文集, vol. 2, Amer. Math. Soc., 1992, p. 79--106.

{[}7{]} \_\_\_\_ 傅里叶积分讲义, Leipzig, 1932, Zbl 0006.11001.

{[}8{]} \_\_\_\_ "酉算子线性族的谱分解", Sitzungsber. Preuß. Akad. Wiss., Phys.-Math. Kl. (1933), p. 371--376, Zbl 0007.16603; 论文集, vol. 1, Amer. Math. Soc., 1992, p. 579--586.

{[}9{]} H. BOHR - "关于几乎周期函数的理论。I. 傅里叶级数理论的推广", Acta Math. 45 (1924), p. 29-127, JfM 50.0196.01.

{[}10{]} J. W. CALKIN - "希尔伯特空间有界算子环中的双边理想与同余关系", Ann. of Math. (2) 42 (1941), p. 839-873, Zbl 0063.00692.

{[}11{]} T. CARLEMAN - 关于具有实对称核的奇异积分方程, Uppsala, 1923, JfM 49.0272.01.

{[}12{]} H. CARTAN and R. GODEMENT -- "局部紧阿贝尔群中的对偶理论与谐波分析", Ann. Sci. Éc. Norm. Supér. (3) 64 (1947), p. 79--99, Zbl 0033.18801; Henri Cartan 文集, vol. III, Springer, 1979, p. 1203--1223.

{[}13{]} H. CARTAN and J-P. SERRE -- "关于紧解析流形的有限性定理", C. R. Acad. Sci., Paris 237 (1953), p. 128--130, Zbl 0050.17701; Henri Cartan 文集, vol. II, Springer, 1979, p. 684--686.

{[}14{]} C. CHEVALLEY - "无限扩张的类域论推广", J. Math. Pures Appl. (9) 15 (1936), p. 359-371, Zbl 0015.15101.

{[}15{]} J. DIEUDONNÉ -- "关于赋范空间的同态", Bull. Sci. Math. (2) 67 (1943), p. 72--84, Zbl 0028.23301; 数学选集, Hermann, 1981, p. 268--281.

{[}16{]} J. DIEUDONNÉ - 泛函分析史, North-Holland Publ. Co., 1981, Zbl 0478.46001.

{[}17{]} P. L. DIRICHLET -- "证明每个首项和公差互素的无限算术级数包含无穷多个素数的定理", Abh. Preuß. Akad. Wiss. 8 (1837), p. 45--81; 文集, vol. 1, Reimer (Berlin), 1889, p. 315--356.

{[}18{]} P. ENFLO - "巴拿赫空间逼近问题的反例", Acta Math. 130 (1973), p. 309-317, Zbl 0267.46012.

{[}19{]} \_\_\_\_ "关于巴拿赫空间不变子空间问题", Acta Math. 158 (1987), no. 3-4, p. 213--313, Zbl 0663.47003.

{[}20{]} J. M. G. FELL - "C*代数的对偶空间", Trans. Amer. Math. Soc. 94 (1960), p. 365-403, Zbl 0090.32803.

{[}21{]} M. FUKAMIYA - "关于盖尔范德和奈马克的一个定理以及 B*代数", Kumamoto J. Sci., Math. 1 (1952), p. 17-22, Zbl 0049.08605.

{[}22{]} C. F. GAUSS - Disquisitiones Arithmeticae, Göttingen, 1801.

{[}23{]} I. M. GELFAND - 论文集, 3 vol., Springer, 1987.

{[}24{]} I. M. GELFAND and M. A. NAIMARK -- "带对合的赋范环及其表示", Izv. Akad. Nauk SSSR, Ser. Mat. 12 (1948), p. 445--480, Zbl 0031.03403; I. M. Gelfand 论文集, vol. I, Springer, 1987, p. 364--400（俄文）。

{[}25{]} R. GODEMENT -- "关于 V. Bargmann 正交关系。II. 一般证明", C. R. Acad. Sci. Paris 225 (1947), p. 657--659, Zbl 0029.19907.

{[}26{]} I. C. GOHBERG -- "关于赋范空间中的线性方程", Dokl. Akad. Nauk SSSR, n. Ser. 76 (1951), p. 477--480, Zbl 0042.11903（俄文）。

{[}27{]} P. R. HALMOS - "谱定理说了什么？", Amer. Math. Monthly 70 (1963), p. 241-247, Zbl 0132.35606; 选集, vol. II, Springer, 1983, p. 59-66.

{[}28{]} W. HEISENBERG -- "关于量子理论形式规则在反常塞曼效应问题中的一种修改", Z. Phys. 26 (1924), p. 291--307, JfM 51.0740.07; 文集：科学原始论文, vol. 1, Springer, 1985, p. 306--325.

{[}29{]} E. HELLINGER and O. TOEPLITZ -- "积分方程与无穷多个未知量的方程", Enzyklopädie der mathematischen Wissenschaften, art. II C 13, Leipzig (Teubner), 1927, JfM 53.0350.01.

{[}30{]} E. HILB - "关于任意函数的积分表示", Math. Ann. 66 (1909), p. 1-66, JfM 39.0380.03.

{[}31{]} D. HILBERT - "线性积分方程一般理论的基础。第一篇通信。", Nachr. Ges. Wiss. Göttingen, Math.-Phys. Kl. (1904), p. 49-91, JfM 35.0378.02.

{[}32{]} \_\_\_\_ "线性积分方程一般理论的基础。第四篇通信。", Nachr. Ges. Wiss. Göttingen, Math.-Phys. Kl. (1906), p. 157--227, JfM 37.0351.03.

{[}33{]} T. H. HILDEBRANDT - "关于完全连续线性变换", Acta Math. 51 (1928), p. 311-318, JfM 54.0427.03.

{[}34{]} E. HILLE - 泛函分析与半群, Colloq. Publ., vol. 31, Amer. Math. Soc., 1948, Zbl 0033.06501.

{[}35{]} M. JAMMER - 量子力学的概念发展, McGraw-Hill, 1966.

{[}36{]} M. KAC - "能听出鼓的形状吗？", Am. Math. Mon. 73 (1966), p. 1-23, Zbl 0139.05603.

{[}37{]} E. R. VAN KAMPEN - "局部双紧阿贝尔群及其特征标群", Ann. of Math. (2) 36 (1935), p. 448-463, Zbl 0011.39401.

{[}38{]} \_\_\_\_ "几乎周期函数与紧群", Ann. of Math. (2) 37 (1936), p. 78--91, Zbl 0014.01702.

{[}39{]} P. LÉVY - "关于傅里叶级数的绝对收敛", Compos. Math. 1 (1934), p. 1-14, Zbl 0008.31201.

{[}40{]} V. I. LOMONOSOV - "与完全连续算子交换的算子族的不变子空间", Funct. Anal. Appl. 7 (1974), p. 213-214, Zbl 0293.47003.

{[}41{]} L. H. LOOMIS - 抽象谐波分析导论, Van Nostrand Co., 1953, Zbl 0052.11701.

{[}42{]} G. W. MACKEY - "局部紧群表示的不可约性。I", Proc. Nat. Acad. Sci. U.S.A. 35 (1949), p. 537-545, Zbl 0035.06901.

{[}43{]} M. MATHIAS - "关于正傅里叶积分", Math. Z. 16 (1923), p. 103--125, JfM 49.0207.03.

{[}44{]} S. MAZUR - "关于线性环", C. R. Acad. Sci., Paris 207 (1938), p. 1025-1027, Zbl 0020.20101.

{[}45{]} J. MERCER - "正型和负型函数及其与积分方程理论的联系", Philos. Trans. R. Soc. Lond., Ser. A 209 (1909), p. 415-446, JfM 40.0408.02.

{[}46{]} S. G. MIKHLIN - "关于希尔伯特空间中线性方程的求解", Dokl. Akad. Nauk SSSR (N. S.) 57 (1947), p. 11-12, Zbl 0029.05101（俄文）。

{[}47{]} F. J. MURRAY and J. VON NEUMANN - "关于算子环", Ann. of Math. (2) 37 (1936), p. 116-229, Zbl 0014.16101; John von Neumann 文集, vol. III, Pergamon Press, 1961, p. 6-120.

{[}48{]} M. NAGUMO - "线性度量环中的一些解析研究", Jpn. J. Math. 13 (1936), p. 61-80, Zbl 3024132.

{[}49{]} J. VON NEUMANN -- "量子力学的数学基础", Göttingen Nachrichten 1927 (1927), p. 1--57, JfM 53.0848.03; 文集, vol. I, Pergamon Press, 1961, p. 151--207.

{[}50{]} \_\_\_\_ "埃尔米特函数算子的一般特征值理论", Math. Ann. 102 (1929), p. 49--131, JfM 55.0824.02; 文集, vol. II, Pergamon Press, 1961, p. 3--85.

{[}51{]} \_\_\_\_ "关于函数运算代数与正规算子理论", Math. Ann. 102 (1929), p. 370--427, JfM 55.0825.02; 文集, vol. II, Pergamon Press, 1961, p. 86--143.

{[}52{]} \_\_\_\_ "群上的几乎周期函数。I", Trans. Amer. Math. Soc. 36 (1934), p. 445--492, Zbl 0009.34902; 文集, vol. II, Pergamon Press, 1961, p. 454--501.

{[}53{]} F. NOETHER -- "关于一类奇异积分方程", Math. Ann. 82 (1920), p. 42--63, JfM 47.0369.02.

{[}54{]} F. PETER and H. WEYL -- "闭连续群的本原表示的完备性", Math. Ann. 97 (1927), p. 737--755, JfM 53.0387.02; Hermann Weyl Gesammelte Abhandlungen, Bd. III, Springer, 1968, p. 58--75.

{[}55{]} L. PONTRYAGIN -- "闭集的一般对偶定理", Verhandlungen Kongreß Zürich, 2, p. 195-197, 1932, JfM 58.0646.08.

{[}56{]} \_\_\_\_ "几乎周期函数与位置分析", C. R. Acad. Sci., Paris 196 (1933), p. 1201--1203, Zbl 0006.42801.

{[}57{]} \_\_\_\_ "闭集的拓扑对偶一般定理", Ann. of Math. (2) 35 (1934), no. 4, p. 904--914, Zbl 0010.18004.

{[}58{]} \_\_\_\_ "拓扑交换群理论", Ann. of Math. (2) 35 (1934), p. 361--388, Zbl 0009.15601.

{[}59{]} F. RIESZ - 无穷多个未知量的线性方程组, Gauthier-Villars, 1913, JfM 44.0401.01.

{[}60{]} \_\_\_\_ "关于线性泛函方程", Acta Math. 41 (1916), p. 71--98, JfM 46.0635.01 ; 全集, t. II, Gauthier-Villars, 1960, p. 1053--1080.

{[}61{]} \_\_\_\_ "关于斯通和博赫纳定理", Acta Litt. Sci. Szeged 6 (1933), p. 184--198, Zbl 0007.30903; 全集, t. II, Gauthier-Villars, 1960, p. 1135--1149.

{[}62{]} J. SCHATZ - "对 M. Fukamiya 所著《关于盖尔范德和奈马克的一个定理以及 B*代数》一文的评论", Math. Reviews 14 (1952), p. 884.

{[}63{]} J. SCHAUDER - "关于线性的完全连续泛函运算", Studia Math. 2 (1930), p. 183-196, JfM 56.0354.01.

{[}64{]} A. SCHUSTER -- "谱的起源", Report on the fifty-second meeting of the British association for the advancement of Science, J. Murray (London), 1882, p. 120--143.

{[}65{]} L. SCHWARTZ - "函数、微分、傅里叶变换概念的推广及其数学和物理应用", Ann. Univ. Grenoble, Sect. Sci. Math. Phys., n. Ser. 21 (1946), p. 57-74, Zbl 0060.27504; 科学全集, t. I, Soc. Math. France, 2011, p. 61-88.

{[}66{]} \_\_\_\_ "同态与完全连续映射", C. R. Acad. Sci., Paris 236 (1953), p. 2472--2473, Zbl 0050.33301 ; 科学全集, t. I, Soc. Math. France, 2011, p. 345--346.

{[}67{]} I. E. SEGAL - "局部紧群的群代数", Trans. Amer. Math. Soc. 61 (1947), p. 69-105, Zbl 0032.02901.

{[}68{]} G. E. SHILOV -- "关于交换赋范环分解为理想直和", Mat. Sb., Nov. Ser. 32 (1953), p. 353--364, Zbl 0052.34203（俄文）。

{[}69{]} S. SOBOLEV - "求解正规双曲型线性方程柯西问题的新方法", Rec. Math. Moscou, n. Ser. 1 (1936), p. 39-71, Zbl 0014.05902.

{[}70{]} S. W. P. STEEN - "算子理论导论：I. 实算子。模", Proc. Lond. Math. Soc. (2) 41 (1936), p. 361-392, Zbl 0014.16201.

{[}71{]} T.-J. STIELTJES - "连分数研究", Ann. Fac. Sci. Toulouse Sci. Math. Sci. Phys. 8 (1894), no. 4, p. J1-J122, JfM 25.0326.01.

{[}72{]} M. H. STONE - 希尔伯特空间中的线性变换及其在分析中的应用, Colloq. Publ., Am. Math. Soc., vol. 15, Amer. Math. Soc., 1932, Zbl 0005.40003.

{[}73{]} \_\_\_\_ "关于希尔伯特空间中的单参数酉群", Ann. of Math. (2) 33 (1932), p. 643--648, Zbl 0005.16403.

{[}74{]} \_\_\_\_ "布尔环理论在一般拓扑中的应用", Trans. Am. Math. Soc. 41 (1937), p. 375--481, Zbl 0017.13502.

{[}75{]} J. T. TATE — “数域中的傅里叶分析与黑克 ζ 函数”, 学位论文, Princeton, 1950; 文集, part I, Amer. Math. Soc., 2016, p. 1–43, Zbl 1407.01030.

{[}76{]} O. TOEPLITZ — “费耶尔定理的代数类似物”, Math. Z. 2 (1918), p. 187–197, JfM 46.0157.02.

{[}77{]} H. WEBER — “证明：每个真正本原二次型都能表示无穷多个素数的定理。”, Math. Ann. 20 (1882), p. 301–330, JfM 14.0141.01.

{[}78{]} \_\_\_\_ “阿贝尔数域理论”, Acta Math. 9 (1887), p. 105–130, JfM 18.0055.04.

{[}79{]} A. WEIL — “关于几乎周期函数”, C. R. Acad. Sci., Paris 200 (1935), p. 38–40, Zbl 0010.35303; 科学全集, vol. I, Springer, 1979, p. 101–103.

{[}80{]} \_\_\_\_ 拓扑群上的积分及其应用, Actualités scientifiques et industrielles, vol. 869, Hermann, 1940, Zbl 0063.08195.

{[}81{]} H. WEYL — “关于差为完全连续的有界二次型”, Rend. Circ. Mat. Palermo 27 (1909), p. 373–392, JfM 40.0395.01; Gesammelte Abhandlungen, Bd. I, Springer, 1968, p. 175–194.

{[}82{]} \_\_\_\_ “关于带奇点的常微分方程及任意函数的相应展开”, Math. Ann. 68 (1910), p. 220–269, JfM 41.0343.01; Gesammelte Abhandlungen, Bd. I, Springer, 1968, p. 248–297.

{[}83{]} \_\_\_\_ “通过线性变换表示连续半单群的理论。I、II、III，补编”, Math. Z. 23 (1925), p. 271–309, JfM 51.0319.01; Gesammelte Abhandlungen, Bd. II, Springer, 1968, p. 543–647.

{[}84{]} \_\_\_\_ “积分方程与几乎周期函数”, Math. Ann. 97 (1926), p. 338–356, JfM 52.0260.02; Gesammelte Abhandlungen, Bd. III, Springer, 1968, p. 38–57.

{[}85{]} \_\_\_\_ “量子力学与群论”, Z. Phys. 46 (1927), p. 1–46, JfM 53.0848.02; Gesammelte Abhandlungen, Bd. III, Springer, 1968, p. 90–135.

{[}86{]} \_\_\_\_ 群论与量子力学, S. Hirzel (Leipzig), 1928, Zbl 0537.22023.

{[}87{]} E. WEYR — “关于双线性型理论”, Monatsh. Math. Phys. 1 (1890), p. 163–236, JfM 22.0141.02.

{[}88{]} N. WIENER — “塔伯尔定理”, Ann. of Math. (2) 33 (1932), p. 1–100, Zbl 0004.05905; 论文选集, MIT Press, 1964, p. 261–360.

{[}89{]} N. WIENER and R. E. A. C. PALEY — “无限阿贝尔群的特征标的解析性质”, Verhandlungen Kongreß Zürich, 2, p. 95, 1932, JfM 58.0445.02.

{[}90{]} E. WIGNER — “关于非齐次洛伦兹群的酉表示”, Ann. of Math. (2) 40 (1939), p. 149–204, Zbl 0020.29601; 文集, part A, Springer, 1993, p. 334–389.

{[}91{]} W. WIRTINGER — “关于黎曼双曲型微分方程积分方法及其在振动问题中应用的贡献”, Math. Ann. 48 (1897), p. 365–389, JfM 27.0281.01.

{[}92{]} S. L. WORONOWICZ — “与 C*代数相关的无界元与非紧量子群”, Comm. Math. Phys. 136 (1991), no. 2, p. 399–432, Zbl 0743.46080.

{[}93{]} B. YOOD — “在线性变换加上完全连续变换后保持的性质”, Duke Math. J. 18 (1951), p. 599–612, Zbl 0043.11901.

\specialchapter{记号索引}

\(\mathcal{L}^{\mathrm{c}}(\mathrm{E};\mathrm{F})\) (紧线性映射) 2 \(\mathrm{E}_{(\mathbf{C})}\) (复化向量空间) 2 \(u_{(\mathbf{C})}\) (复化线性映射) 3 \(\mathcal{L}^{\mathrm{f}}(\mathrm{E};\mathrm{F})\) (有限秩线性映射) 4 \(\mathcal{C}_0(\mathrm{X};\mathrm{K}), \mathcal{C}_0(\mathrm{X})\) (在无穷远处趋于0的连续函数) 20 \(\mathcal{N}^{p,q}(\mathrm{X} \times \mathrm{Y})\) (核空间) 27 \(u_k\) (带核的积分算子) 29 \(k_\circ\)（映射 \(y \mapsto k_y\) ，其中 \(k_y(x) = k(x,y)\) ）30 \(\mathcal{L}^{p',q'}(\mathrm{X},\mathrm{Y},\mu,\nu)\) (核空间) 31 \(u \equiv v\) (模 \(\mathcal{L}^{\mathrm{f}}(\mathrm{E};\mathrm{F})\) 的同余) 40 \(\mathcal{F}(\mathrm{E};\mathrm{F})\) (弗雷德霍姆映射) 41\\
ind(u) (弗雷德霍姆映射的指标) 43 \(\mathcal{I}(\mathrm{E};\mathrm{F})\) (E 到 F 的同构) 57 \(\mathcal{M}\mathcal{D}(\mathrm{E};\mathrm{F})\) (单射直态射) 57 \(\mathcal{E}\mathcal{D}(\mathrm{E};\mathrm{F})\) (满射直态射) 57 \(\mathcal{R}(\mathrm{E})\) (里斯自同态) 60\\
((u)) (u 的余范数) 61 \(\mathcal{M}(\mathrm{E};\mathrm{F})\) (严格单射态射) 67 \(\mathcal{Q}\mathcal{M}(\mathrm{E};\mathrm{F})\) (核有限维的严格态射) 67 \(\mathcal{E}(\mathrm{E};\mathrm{F})\) (连续满射态射) 70 \(\mathcal{Q}\mathcal{E}(\mathrm{E};\mathrm{F})\) (像具有有限余维的态射) 70 \(\mathcal{C}alk(\mathrm{E})\) (巴拿赫空间 E 的卡尔金代数) 75 \(\mathrm{N}_{\lambda}\) (幂零空间) 78 \(\mathrm{I}_{\lambda}\) (余幂零空间) 78 \(e_{\lambda}(u)\) (谱投影) 82 \(m_{\lambda}(u)\) (谱重数) 82 \(\mathrm{Sp}_{\mathrm{s}}(u)\) (谱的敏感点) 83 \(\mathrm{Sp}_{\mathrm{e}}(u)\) (本质谱) 83 \(\mathrm{Sp}_{n}(u)\) (谱中使得 \(u - \lambda 1_{\mathrm{E}}\) 具有指标 n 的元素) 86 \(\mathrm{Sp}_{\omega}(u)\) (\(\mathrm{Sp}_{n}(u)\) 的并集的补集) 86 \(\mathcal{D}_{\mathrm{B}}(\mathrm{E})\) (基 B 中对角的自同态) 146 \(\mathrm{I}_{\mathrm{E}}\) (F ≠ E 的有限维子空间的维数) 151 \(\lambda_n(u)\) (u 的特征值递减序列) 152

\(r_{\mathrm{F}}(u), \mathrm{R}_{\mathrm{F}}(u)\) （瑞利商）\ldots.. 153 \(\mathcal{F}_n\) （维数 \(n\) 的子空间）\ldots.. 153 \((\alpha_n(u))_{n \in \mathrm{J}}\) （奇异值序列）\ldots.. 156 \((\alpha_n(u))_{n \in \mathrm{I_E}}\) （奇异值的扩展序列）\ldots.. 156 \(\mathcal{L}_1(\mathrm{E})\) （有限迹自同态）\ldots.. 164 \(\|u\|_1\) （\(\mathcal{L}_1(\mathrm{E};\mathrm{F})\) 上的范数）\ldots.. 167 \(\mathcal{L}_1(\mathrm{E};\mathrm{F})\) （核型线性映射）\ldots.. 169 \(\mathcal{L}_{\mathrm{u}}(\mathrm{X})\) （普遍可测函数）\ldots.. 183 \(\mathcal{L}_{\mathrm{u}}^{\infty}(\mathrm{X})\) （有界普遍可测函数）\ldots.. 183 \(m_g, \widetilde{m}_g\) （乘法自同态）\ldots.. 186 \(\mathcal{G}_{\mathrm{A}}\) （盖尔范德变换）\ldots.. 190 \(\mathcal{H}_{\mathrm{A}}\) （盖尔范德变换的逆）\ldots.. 190 \(\mathcal{F}\) （傅里叶变换）\ldots.. 196 \(\overline{\mathcal{F}}\) （余傅里叶变换）\ldots.. 196 \(\mathcal{D}(\mathrm{U})\) （U 上的测试函数）\ldots.. 200 \(\mathcal{D}_{\mathbf{R}}(\mathrm{U})\) （实值测试函数）\ldots.. 201 \(\mathcal{D}'(\mathrm{U})\) （U 上的广义函数）\ldots.. 203 \(\partial^\alpha f\) （广义函数的偏导数）\ldots.. 208 \(\mathcal{S}(\mathbf{R}^n)\) （\(\mathbf{R}^n\) 上的施瓦茨函数）\ldots.. 209 \(\mathcal{S}'(\mathbf{R}^n)\) （\(\mathbf{R}^n\) 上的缓增广义函数）\ldots.. 215 \(\mathcal{M}^t(\mathbf{R}^n)\) （缓增测度空间）\ldots.. 216 \(W^{k,p}(U)\) （指标 \(k\) 、指数 \(p\) 的索伯列夫空间）\ldots.. 222 \(W_0^{k,p}(U)\) （\(\mathcal{D}(U)\) 在 \(W^{k,p}(U)\) 中的闭包）\ldots.. 222 \(H^k(U)\) （指标 \(k\) 、指数2的索伯列夫空间）\ldots.. 222 \(H_0^k(U)\) （\(\mathcal{D}(U)\) 在 \(H^k(U)\) 中的闭包）\ldots.. 222 \(\mathcal{P}(\mathrm{E};\mathrm{F})\) （从 E 到 F 的部分算子）\ldots.. 224 \(\mathcal{P}(\mathrm{E})\) （E 上的部分算子）\ldots.. 224 dom( \(u\) ) （\(u\) 的定义域）\ldots.. 224 \(1_{\mathrm{D}}\) （定义域为 D 的恒等部分算子）\ldots.. 225 \(u \circ v\) （部分算子 \(u\) 与 \(v\) 的复合）\ldots.. 225 \(u \subset v\) （\(v\) 是 \(u\) 的扩张）\ldots.. 226 \((u_i)_{i \in I}\) （乘积部分算子）\ldots.. 226 \(u + v\) （部分算子的和）\(u\)、\(v\)\ldots.. 226 \(\overline{u}\) （\(u\) 的闭包）\ldots.. 228 \(E_u\) （由 \(u\) 的定义域定义的赋范空间）\ldots.. 230 \(\|x\|_u\) （空间 \(E_u\) 上的范数）\ldots.. 230 \((x | y)_u\) （定义希尔伯特空间 \(E_u\) 的内积）\ldots.. 230 \(m_g, \widetilde{m}_g\) （乘法部分算子）\ldots.. 231 \(u^*\) （\(u\) 的伴随）\ldots.. 235 \(\mathcal{A}(\mathrm{E})\) （自伴部分算子）\ldots.. 238 \(^t\mathrm{D}\) （转置微分算子）\ldots.. 242 \(H_D\) （微分算子 \(D_+\) 的定义域）\ldots.. 242 \(D_-\) （定义域为 \(\mathcal{D}(U)\) 的微分算子）\ldots.. 242 \(D_+\) （定义域为 \(H_D\) 的微分算子）\ldots.. 242 \(R(u,\lambda)\) （\(u\) 的预解式）\ldots.. 244 \(\mathrm{PSp}_{\varepsilon}(u)\) （\(\varepsilon\)-伪谱）\ldots.. 250 \(b(u)\) （\(u\) 的有界化）\ldots.. 262 \(\Omega(\mathrm{E})\) （使 \(v\) 正且单射的 \(1_{\mathrm{E}} - v^*v\) 的集合）\ldots.. 264 \(B(v)\) （部分算子 \(v\circ((1_{\mathrm{E}} - v^*v)^{1/2})^{-1}\) ）\ldots.. 264

\(\beta (z) = z / \sqrt{1 + |z|^2}\) 266 \(D_f\) （\(f(u)\) 的定义域）271 \(f(u)\) （普遍可测函数演算）272 \(p_A\) （由 A 定义的谱投影）277 \(u^\alpha\) （函数演算）289 \(|u|\) （\(u\) 的绝对值）290 \(q_u\) （与 \(u\) 关联的正部分形式）291 \(E_q\) （预希尔伯特空间）292 \((x|y)_q\) （\(E_q\) 上的内积）292 \(\| x \| _q\) （\(E_q\) 上的范数）292 \(u_q\) （表示 \(q\) 的部分算子）292 \(\mathrm{Sp}_{\mathrm{s}}(u)\) （\(u\) 的敏感谱）297 \(\mathrm{Sp}_{\mathrm{e}}(u)\) （\(u\) 的本质谱）297 \(\mathrm{Sp_h}(u)\) （谱的上部）298 \(\varrho_{e}\) （本质谱的下界）298 \(\mathrm{Sp_b}(u)\) （谱的下部）298 \(\mathcal{F}_n\) （维数 \(n\) 的子空间）299 \(\mathcal{F}_n^u\) （\(\mathrm{dom}(u)\) 中维数 \(n\) 的子空间）299 \(\widetilde{r}_{\mathrm{F}}(u), \widetilde{\mathrm{R}}_{\mathrm{F}}(u)\) （瑞利商）300 \(\|v\|_u\) （关于 \(v\) 的 \(u\) 的相对范数）306\\
vp （柯西主值）337 \(\chi_\varrho\) （表示的特征标）374 \(\mathrm{Res}_{H}^{G}(\varrho)\) （\(\varrho\) 到 H 的限制）375 \(\mathrm{Hom}_{G}(\varrho_1,\varrho_2)\) （G-态射空间）375 \(\varrho^n = n \varrho\) （表示 \(\varrho\) 的副本之和）376 \(\check{\varrho}\) （逆步表示）377 \(\mathrm{Sesq_G}(\varrho, \pi), \mathrm{Sesq_G(E,F)}\) （G-不变半双线性形式）380 \(\overline{\varrho}\) （共轭表示）382 \(\varrho_1 \boxtimes \varrho_2\) （外张量积）385 \(\varrho_1 \otimes \varrho_2\) （张量积）385 \(\varrho^{\otimes n}\) （\(\varrho\) 的张量幂）385 \(\Upsilon(G)\) （矩阵系数）386 \(\Theta(G)\) （有限维矩阵系数）386 \(\operatorname{cl}(\pi)\) （酉表示的类）392 \(\widehat{G}\) （酉对偶）393 \(M_\pi (\varrho)\) （同型分量）394 \(\mathcal{M}_c(G)\) （紧支撑测度）400 \(\mathcal{M}^1(G)\) （有界测度）401 \(\gamma_{G}^{(p)}\) （左正则表示）405 \(\delta_{G}^{(p)}\) （右正则表示）405 \(\varrho_{G}^{(p)}\) （双正则表示）406 \(\gamma_{G}\) （\(L^2\) 中的左正则表示）407 \(\delta_{G}\) （\(L^2\) 中的右正则表示）407 \(\varrho_{G}\) （\(L^2\) 中的双正则表示）407 \(\mathcal{F}_{\pi}(G)\) （π-等变函数）408 \(\mathcal{K}_{\pi}(G)\) （在）\(\mathcal{F}_{\pi}(G)\) 中模 H 具有紧支集的连续函数）\ldots{} 408 \(N_p(f)\) （在）\(\mathcal{F}_π(G)\) 上的半范数）\ldots{} 408 \(\mathcal{L}_{\pi}^{p}(G)\) （模 H 对幂 \(p^e\) 可积的 π-等变函数）\ldots{} 409

\(L_{\pi}^{p}(G)\) （与 \(\mathcal{L}_{\pi}^{p}(G)\) 关联的巴拿赫空间）\ldots.. 410 \(\mathcal{L}_{\pi}^{2}(G)\) （模 H 平方可积的 \(\pi\)-等变函数）\ldots.. 415 \(L_{\pi}^{2}(G)\) （与 \(\mathcal{L}_{\pi}^{2}(G)\) 关联的希尔伯特空间）\ldots.. 415 \(\operatorname{Ind}_{H}^{G}(\pi)\) （由 H 的表示 π 诱导的 G 表示）\ldots.. 416 \(\gamma_{G,\chi}\) （在 \(\mathcal{F}_{\chi}(G)\) 或 \(L_{\chi}^{2}(G)\) 上的左正则表示）\ldots.. 417 \(\delta_{G,\chi}\) （在 \(\mathcal{F}_{\chi}(G)\) 或 \(L_{\chi}^{2}(G)\) 上的右正则表示）\ldots.. 417 \(\gamma_{G,\chi}^{(p)}\) （\(L_{\chi}^{p}(G)\) 中的左正则表示）\ldots.. 417 \(\delta_{G,\chi}^{(p)}\) （\(L_{\chi}^{p}(G)\) 中的右正则表示）\ldots.. 417 \(\gamma_{G,\chi}\) （\(L_{\chi}^{2}(G)\) 中的酉左正则表示）\ldots.. 417 \(\delta_{G,\chi}\) （\(L_{\chi}^{2}(G)\) 中的酉右正则表示）\ldots.. 417 \(\varrho_{G,\chi}\) （\(L_{\chi}^{2}(G)\) 中的酉双正则表示）\ldots.. 417 \(c_{Z}(\pi)\) （形式次数）\ldots.. 423 \(\mathrm{Noy}_{+}(\mathrm{X})\) （普遍正核）\ldots.. 434 \(A_{+}^{\prime}\) （连续正线性泛函）\ldots.. 436 \(\mathrm{Pos}(G)\) （正定函数）\ldots.. 443 \(\mathrm{Pos}_{1}(G)\) （满足 \(\varphi(e)=1\) 的正定函数）\ldots.. 443 \(\varrho_{G}(\pi)\) （\((\overline{\pi}\boxtimes\pi)\)-同型分量的 \(\varrho_{G}\)）\ldots.. 462 \(\Theta(G^{\sharp})\)，\(Z\Theta(G)\) （中心矩阵系数）\ldots.. 467 \(R(G)\) （格罗滕迪克环）\ldots.. 469 \(\mathrm{F}(\widehat{G})\) （End(Eπ) 的乘积代数）\ldots.. 470 \(\mathrm{F}_{b}(\widehat{G})\) （有界族的子代数）\ldots.. 470 \(\|u\|_{\mathrm{F}_{b}}\) （\(\mathrm{F}_{b}(\widehat{G})\) 中的范数）\ldots.. 470 \(\mathrm{F}_{0}(\widehat{G})\) （趋于0的族的子代数）\ldots.. 470 \(\overline{\mathcal{F}}_{G}(\nu)\) （ν 的傅里叶余变换）\ldots.. 471 \(\overline{\mathcal{F}}_{G}(f)\) （f 的傅里叶余变换）\ldots.. 471 \(\|u\|_{2}\) （End(Eπ) 中的范数）\ldots.. 471 \(\mathrm{F}^{2}(\widehat{G})\) （End(Eπ) 的希尔伯特直和）\ldots.. 472 \(\mathcal{F}_{G}(f)\) （f 的傅里叶变换）\ldots.. 473 \(\mathcal{C}(\widehat{G})\) （\(\widehat{G}\) 上的函数）\ldots.. 475 \(M^{1}(G^{\sharp})\) （中心有界测度）\ldots.. 475 \(\mathrm{FS}(\pi)\) （弗罗贝尼乌斯--舒尔指标）\ldots.. 476 \(M_{2k}(\varrho)\) （阶为 2k 的绝对矩）\ldots.. 476 \(S^{2}(E)\) （A-模的对称平方）\ldots.. 478 \(\wedge^{2}E\) （A-模的外平方）\ldots.. 478

\specialchapter{术语索引}

A 卡尔金代数，75 自然的，130 弗雷德霍姆择一定理，79 拉森的，476 弱混合映射，
322
紧线性映射，2 希尔伯特--施密特，23 有限秩的，4 有限迹的，164 由核定义的，25 核型的，169 严格的，2 逼近空间的---，14 逼近性质，14 自伴部分等距---，238 本质部分等距---，238 自守（表示---），405

B 基底（自然代数---），131 底部（谱的---部分），298 双交换子（---定理），326 双正则（表示---），406 博赫纳（---定理），455

有界（表示---），374 有界化，262 单位球，1 布里顿（---引理），506 C 函数演算，185 普遍可测的，272 卡尔金（---代数），75 特征标，374 中心的，390 亏指标偶的---，256 ---的亏指标，256 外平方，478 对称的，478 柯西（---测度），515 柯西--柯瓦列夫斯卡娅（---定理），343 中心特征标---，390 中心测度---，402 矩阵系数，386 有限维的，386 对角的，386 部分算子的核心，231 紧（线性映射---），2 完全连续，104，107 同型分量，394 边界条件，258 余幂零空间，45

共轭（表示---），382 余范数，61 切格常数，323 盖尔范德--奈马克--塞加尔构造，433 逆步（表示---），377 广义函数与函数的卷积，339 对应，224 余傅里叶变换，196，471 分布的，219 余傅里叶变换，471 盖尔范德对，512 卡日丹的，499 亏指标偶，256 外尔等分布判据，470 循环希尔伯特实现---，433，438，443 向量---，191

D 极分解，291 形式次数，423 广义函数的偏导数，208 对角自同态---，24 基中的对角自同态---，146 基中的对角部分算子---，234 直态射---，56 直向量子空间---，55 广义函数，203 紧支撑的，336 缓增的，215 部分算子的定义域，224 酉对偶，393

E 弗雷德霍姆自同态，41 乘法的，186 里斯的，46

有限迹的，23，164 对角的，24 基 B 中对角的，146 卡日丹集，499 预解式，244 薛定谔方程，431 积分的，80 等变（函数---），408 逼近空间，14 施瓦茨空间，209 索伯列夫空间，222 本质的（谱---），297 自伴扩张，255 弗里德里希斯的，295 几乎解析的，281 部分算子的扩张，226 F 层，204 部分算子的闭包，228 向量丛，36 忠实（表示---），374 多项式增长函数，214 中心的，402 勒夫纳的，326 施瓦茨的，209 负定型的，495 正定型的，442 插值的，327 代表的，386 球的，513 测试的，200 部分埃尔米特形式，291 正线性泛函，436 部分正的，291 关联的部分正的，291 闭的部分正的，292 不变半双线性的，380 形式的（---次数），423 形式对称（---微分算子），243 埃尔弗--舍斯特兰德公式，284 莱布尼茨公式，196，209

普朗谢雷尔的，474 斯通的，353 迹公式，488 弗雷德霍姆择一定理的---，79 ---的映射，41 ---的自同态，41 ---的映射的指标，43 弗里德里希斯（---扩张），295 弗罗贝尼乌斯--舒尔（---指标），476 G 盖尔范德--奈马克（---定理），442 盖尔范德--奈马克--塞加尔（---构造），433 盖尔范德--赖科夫（---定理），454 无穷小生成元，428 部分算子的图像，224 H 汉堡（---矩问题），359 高（谱的---部分），298 埃尔弗--舍斯特兰德（---公式），284 海林格--特普利茨（---定理），240 希尔伯特实现---，433，438，443 循环希尔伯特实现---，433，438，443 希尔--吉田（---定理），346 I 恒等（对合自构---），478 μ-本质像，181 弗罗贝尼乌斯--舒尔指标，476 指标，43 亏指标的，256 加藤不等式，344 不可约矩阵---，132 不可约表示---，378 表示的同构，375 同型的（分量---），394 K 加藤

---的不等式，344 加藤--雷利希（---定理），306 ---的卡日丹对，499 ---的卡日丹集，499 ---的卡日丹性质 {]}T{[}，498 凯斯滕（---定理），323 克赖因--鲁特曼（---定理），94

L 拉普拉斯算子，243 狄利克雷的，310 拉森（---择一定理），476 莱布尼茨（---公式），196 舒尔引理，12，386 卢伊（---定理），343 勒夫纳（---定理），330 外尔定律，311 洛莫诺索夫（---定理），13

M 不可约矩阵，132 本原的，138 可约的，132 极大–极小，154，300 默瑟（---定理），177 中心测度，402 柯西的，515 谱的，191，268 缓增的，216 极小–极大，154，300 下有界（部分算子---），238 绝对矩，476 测度的矩，359 ---的矩问题（汉堡的），359 ---的矩问题（斯蒂尔切斯的），364 穆尔--彭罗斯（---伪逆），320 表示的态射，375 直的，56 ---的乘法自同态，186 ---的乘法部分算子，231 重数，245 ρ 中的 π 的，398

谱的，82 有限谱的，82

N 自然（代数---），130 负的（负定型函数---），495 模 H 零测（集合---），
408
幂零空间，45，48 非退化（表示---），
373
正规（部分算子---），265 相对范数，306 部分算子的核，225 普遍正的，433 核型（线性映射---），
169
\lettergroup{0}

带指标的算子，41 带核的，25 紧的，2 弗雷德霍姆的，41 希尔伯特--施密特的，23 马尔可夫的，322 乘法的，186 有限秩的，4 里斯的，46 斯图姆--刘维尔的，350 有限迹的，23 沃尔泰拉的，118 微分的，235 标量微分的，235 转置微分的，242 带核的积分的，29

部分算子，224 定义域稠密的，227 具有紧预解式的，309 伴随的，235 自伴的，238 双射的，225 ---的核心，231 乘法的，231 ---的定义域，224 本质自伴的，238 ---的扩张，226 可闭的，227 闭的，227

---的闭包，228 ---的图像，224 在---下的像，225 在---下的原像，225 单射的，225 正规的，265 ---的核，225 正的，238 乘积的，226 逆的，226 ---的约化，226 相对有界的，306 表示正形式的，
292
---的预解式，244 和，226 ---的谱，244 满射的，225 对称的，238 正交性（---的关系），424，
458
\lettergroup{P}

相差局部零测集意义下相等的集合，425 佩龙--弗罗贝尼乌斯（---定理），

彼得--外尔（---定理），462 普朗谢雷尔（---公式），474 谱的敏感点，83 第一类正交多项式，361 第二类的，363 ---型的正定函数，442 部分算子---，238 正的（线性泛函---），436 本原的（矩阵---），138 极大–极小原理，154，300 极小–极大原理，154，300 酉表示的张量积，385 酉表示的外张量积，385 谱投影，54，277 逼近性质，14 卡日丹的 {]}T{[}，498

模 H 几乎处处为真，408 伪谱，250 Q 拟逆，40 典范的，48 拟同构，41 R 希尔伯特实现，433，438，443 循环希尔伯特实现，433，438，443 可约（矩阵---），132 约化，226 正则表示---右，405 表示---左，405 正交关系，424，458 Γ-自守表示，405 双正则的，406 有界的，374 共轭的，382 逆步的，377 平方可积的，420 模中心平方可积的，420 正交型的，476 四元数型的，476 实型的，476 辛型的，476 忠实的，374 诱导的，416 不可约的，378 ---的表示同构，375 线性的，374 连续线性的，374 ---的态射，375 非退化的，373 酉表示的张量积，385 张量幂，385 拟正则的，486 商的，376 右正则的，405 左正则的，405 ---对子群的限制，375 ---的副本之和，376

---的希尔伯特直和，酉，384 子---，375 平凡的，374 酉的，379 预解式，244 具有紧---的部分算子，309 表示的限制，375 S 绍德尔（---定理），9 薛定谔（---方程），431 舒尔（---引理），12，386 施瓦茨空间的---，209 ---函数，209 连续半群，345 半单（酉表示---），391 敏感的（谱的---点），83，297 斯梅尔（---定理），124 索伯列夫（---空间），222 分部求和，160 适应于 u 的子空间，153，299 坐标的，132 不变的，12 直向量，55 子表示，375 由一个子集生成的，376 谱测度---，268 ---的重数，82 有限重数---，82 ---投影，277 ---定理，194，266 谱，244 复的，54 连续的，245 本质的，83，297 ---的敏感点，83 剩余的，245 敏感的，297 稳定的，86 斯蒂尔切斯（---矩问题），364 斯通的---公式，353

---定理，428 特征值序列，152 广义函数的支撑，336 符号，127 对称（部分算子---），238

\lettergroup{T}

博赫纳定理，455 博苏克--乌拉姆定理，64 柯西--柯瓦列夫斯卡娅定理，343 盖尔范德--奈马克定理，442 盖尔范德--赖科夫定理，454 海林格--特普利茨定理，240 希尔--吉田定理，346 加藤--雷利希定理，306 凯斯滕定理，323 克赖因、克拉索塞利斯基和米尔曼定理，64 克赖因--鲁特曼定理，94 卢伊定理，343 勒夫纳定理，330 洛莫诺索夫定理，13 默瑟定理，177 佩龙--弗罗贝尼乌斯定理，101 彼得--外尔定理，462 椭圆正则性定理，341 雷利希--孔德拉绍夫定理，335 绍德尔定理，9 斯梅尔定理，124 斯通定理，428 冯·诺伊曼定理，258 双交换子定理，326 施瓦茨核定理，342 谱定理，194，266 特普利茨（---算子），127 弱拓扑，6 C \(^{r}\)-紧收敛的，34 弱的，449 T-预解式，285 迹

有限迹映射，164 迹公式，488

傅里叶变换，196，473 盖尔范德变换，190

傅里叶变换，473 T-预解式（拓扑---），285 平凡（表示---），374

U

自然含单位代数

酉对偶---，393 ---表示的张量积，385 表示---，379 半单表示---，391 表示---的希尔伯特直和，384

逼近单位元

普遍可测函数演算---，272 函数---，183

V

柯西主值，337 特征值，245 特征值（---族），146，234
234
奇异值，151，156 循环向量，191，376 G-有限，376 沃尔泰拉（---算子），118 冯·诺伊曼的---定理，258 ---的双交换子定理，
326
\lettergroup{W}

外尔等分布判据---，470 ---不等式，162 ---定律，311

\specialchapter{自伴性判据}

这里汇总正文中陈述的、可用来证明某些部分算子自伴的判据。

固定一个复希尔伯特空间 E 以及 E 上的部分算子 u。在下列每一种情形中，该算子都是自伴的：

\begin{enumerate}
\def\labelenumi{(\roman{enumi})}
\item
  存在一个单射且为埃尔米特的 \(v \in \mathcal{L}(\mathrm{E})\) ，使得 \(u = v^{-1}\) 。（IV，第239页命题9）；
\item
  算子 \(u\) 是对称的，且 \(\mathrm{dom}(u) = \mathrm{E}\) （IV，第240页命题10）；
\item
  算子 \(u\) 是对称的，并诱导从 \(\mathrm{dom}(u)\) 到 E 的双射（IV，第240页命题10的推论）；
\item
  存在 \(\lambda \in \mathbf{C}\) ，使得 \(u + \lambda 1_{\mathrm{E}}\) 和 \(u + \overline{\lambda} 1_{\mathrm{E}}\) 均为满射（IV，第240页命题11）；
\item
  存在一个定义域稠密的闭部分算子 \(v\) ，使得 \(u = v^{*} \circ v\) （IV，第241页命题12）；
\item
  算子 \(u\) 是对称且闭的，并且其亏指标为 \((0,0)\) （IV，第257页命题26的推论），也就是说，\(\mathrm{Ker}(u^{*} - i) = \mathrm{Ker}(u^{*} + i) = \{0\}\) ；
\item
  空间 E 具有 \(\mathrm{L}^{2}(\mathrm{X},\mu)\) 的形式，其中 X 是局部紧拓扑空间，\(\mu\) 是 X 上的测度，而 u 是乘法算子，即乘以一个在 X 上 \(\mu\) -可测且局部地 \(\mu\) -几乎处处取实值的函数（命题23的推论）；
\item
  存在 E 上的正规部分算子 v，以及一个取实值的普遍可测函数 \(f \in \mathcal{L}_{\mathrm{u}}(\mathrm{Sp}(v))\) ，使得 \(u = f(v)\) （IV，第273页推论1）；
\item
  存在 E 上的部分埃尔米特形式 q，使 u 是表示 q 的部分算子（IV，第293页命题14）；
\item
  有 \(u = v + w\) ，其中 v 是自伴部分算子，w 是关于 v 相对有界且满足 \(\|w\|_{v} < 1\) 的对称部分算子（IV，第306页定理5）；
\item
  算子 \(u\) 是 E 上的酉表示 \(\varrho\) （其群为 \(\mathbf{R}\)）的无穷小生成元（V，第428页定理1）。
\end{enumerate}
